% 连通分量：三角形加孤立点；单向链1->2->3的强、弱分量；有向环的单一强分量。
\begin{tikzpicture}[panelmark/.style={font=\figurefont\fontsize{8}{10}\selectfont,text=FigureInk,anchor=north,inner sep=0pt,outer sep=0pt},
>=Stealth,
  v/.style={circle,draw=FigureStroke,fill=FigureBlueFill,line width=.8pt,minimum size=6.5mm,inner sep=0pt},
  e/.style={draw=FigureStroke,line width=.8pt},
  a/.style={draw=FigureStroke,line width=.8pt,->,shorten >=1pt},
  grp/.style={draw=FigureBlue,fill=FigureBlueFill,fill opacity=0,rounded corners=6pt,line width=.7pt},
  sgrp/.style={draw=FigureGreen,dashed,line width=.7pt,rounded corners=6pt},
  note/.style={font=\FigureNoteFont\fontsize{7.5}{9}\selectfont,text=FigureMuted,align=center}]
\begin{scope}[local bounding box=panelrow]
\begin{scope}[local bounding box=panela]
  \node[v] (1) at (0,.7) {$1$}; \node[v] (2) at (-.75,-.5) {$2$}; \node[v] (3) at (.75,-.5) {$3$};
  \node[v] (4) at (1.9,.1) {$4$};
  \draw[e] (1)--(2)--(3)--(1);
  \draw[grp] (-1.25,-.95) rectangle (1.25,1.15);
  \draw[grp] (1.5,-.3) rectangle (2.3,.5);

\end{scope}
\begin{scope}[xshift=4.5cm,local bounding box=panelb]
  \node[v] (1) at (0,0) {$1$}; \node[v] (2) at (1.4,0) {$2$}; \node[v] (3) at (2.8,0) {$3$};
  \draw[a] (1)--(2); \draw[a] (2)--(3);
  \foreach \p in {1,2,3}{\draw[sgrp] ($(\p)+(-.45,-.45)$) rectangle ($(\p)+(.45,.45)$);}
  \draw[grp] (-.65,-.65) rectangle (3.45,.65);

\end{scope}
\begin{scope}[xshift=10.4cm,local bounding box=panelc]
  \node[v] (1) at (0,.7) {$1$}; \node[v] (2) at (-.75,-.5) {$2$}; \node[v] (3) at (.75,-.5) {$3$};
  \draw[a] (1)--(2); \draw[a] (2)--(3); \draw[a] (3)--(1);
  \draw[sgrp] (-1.25,-.95) rectangle (1.25,1.15);

\end{scope}
\end{scope}
% Subfigure marks: 8 pt, centred below each panel; common top 3 mm below this row.
\coordinate (panelrowbase) at ($(panelrow.south)+(0,-3mm)$);
\node[panelmark] at (panela.center |- panelrowbase) {(a)};
\node[panelmark] at (panelb.center |- panelrowbase) {(b)};
\node[panelmark] at (panelc.center |- panelrowbase) {(c)};
\end{tikzpicture}
